\chapter{Thank You}\label{cha:thanks}\markboth{Thank You}{Thank You}
% Every entry starts at the margin: a name half-indented reads as an
% accident, and in a credits list the eye wants one left edge.
\begingroup
\setlength{\parindent}{0pt}
\setlength{\parskip}{5pt plus 1.5pt minus 1pt}

The K4510 is a machine that never existed. Almost everything in it
that \emph{does} exist was written by somebody else first, and given
away. This chapter is the thanks; the Licences chapter (page~\pageref{cha:licences})
is the paperwork.

\section{The heart of the machine}
\textbf{Gábor Lénárt} (LGB) wrote the 45GS02 CPU core in
\href{https://github.com/lgblgblgb/xemu}{Xemu}, the MEGA65 emulator.
It runs here unchanged --- not a line altered, not a bug fixed --- and
every instruction this book describes is his code executing. There
would be no K4510 without it.

\textbf{Dag Lem} wrote reSID, the SID emulation that has been the
reference for thirty years. It was this machine's voice for its whole
first summer --- four 6581s, as shipped in VICE~3.3 --- and it is not
here any more (Appendix~\ref{cha:sound} says why, and the reason is not
that anything was wrong with it). The thanks stands as it was.

\textbf{Jarek Burczyński} and \textbf{Tatsuyuki Satoh} wrote
\texttt{fmopl}, the OPL2 emulation in MAME that this machine's YM3812 is
--- by way of VICE, vendored unaltered. Nine voices of FM that any
AdLib register list from 1987 can drive is their work, not ours.

\textbf{The VICE team}, whose decades of emulation scholarship ---
documentation, test programs, arguments settled in code --- this
project consulted at nearly every turn.

\section{The tongues}
\textbf{Lee Davison} (1966--2013) wrote EhBASIC, the machine's first
BASIC and still the one that comes up in ROM. He is not here to be
asked about the graphics and sound words bolted onto it; the machine
carries his name in its \verb|README| and its startup banner instead.

\textbf{Richard T. Russell} wrote BBC BASIC, and has kept writing it
for forty years. The console edition of his BBC BASIC for SDL~2.0
(BBCTTY) is what runs on the Tube co-processor in
Chapter~\ref{cha:tube}, under the zlib licence he publishes it with.
``BBC BASIC'' is the name of his interpreter; this book uses it only to
say what is running, and claims nothing in it.

\textbf{Marcelo Dantas} (``Mockba the Borg'') wrote
\href{https://github.com/MockbaTheBorg/RunCPM}{RunCPM}, which is the
whole of Chapter~\ref{cha:cpm}: a complete CP/M~2.2 with its own CCP,
vendored here unmodified and simply handed a Z80's worth of address
space.

\textbf{Scot W. Stevenson}, \textbf{Sam Colwell} and \textbf{Patrick
Surry} wrote \href{https://github.com/SamCoVT/TaliForth2}{Tali
Forth~2} and put it in the public domain --- a Forth written to be
read, which is why Chapter~\ref{cha:forth} can be honest about how it
works.

\textbf{Tomasz Biela} (``tebe'') wrote
\href{https://github.com/tebe6502/Mad-Pascal}{Mad Pascal} and
\href{https://github.com/tebe6502/Mad-Assembler}{MADS}; the K4510
target in \verb|pascal/| is a guest in his compiler.
\textbf{Wojciech Bociański} (``bocianu''), whose Neo6502 target showed
how such a guest should behave.

\textbf{Steve Wozniak} wrote Wozmon in 1976 in 256 bytes, and the
monitor in this machine is still recognisably it.

\textbf{Jim Butterfield} (1936--2007) wrote Supermon, and gave it away
as he gave away everything: the machine language monitor a generation of
Commodore programmers learned on, published in \emph{Compute!} for
anyone to type in. Supermon+64 V1.2 is \verb|SUPERMON| here, ported and
not rewritten. \textbf{J.\,B. Langston} restored and commented the
source (\href{https://github.com/jblang/supermon64}{jblang/supermon64}),
which is what made porting it a day's work rather than a month's; he asks
only for the attribution, and it is given gladly.

\textbf{Michael Steil} maintains
\href{https://github.com/mist64/msbasic}{msbasic}, and
\textbf{Microsoft} released 6502 BASIC~1.1 under the MIT licence in
2025 --- between them, the machine's next native BASIC.

\section{Letters on the screen}
Fonts are the part of a computer you look at longest, and clean ones
with a clear pedigree are hard to come by.

\textbf{The Linux kernel} contributors --- the 8x8 console font
(\verb|font_8x8.c|) that got the text mode on its feet and is still the
default.

\textbf{Paul Gardner-Stephen} and \textbf{Roman Standzikowski}
(FeralChild64) --- the clean-room C64/C65 chargen from
\href{https://github.com/MEGA65/open-roms}{MEGA65 open-roms}, and
\textbf{Retrofan}, whose PXLfont travels with it by permission.

\textbf{Ville-Matias Heikkilä} (``Viznut'') --- unscii, placed in the
public domain.

\textbf{Damian Vila} --- \href{https://codeberg.org/Dmian/font-bescii}{BESCII},
the PETSCII spirit with a clean pedigree, CC0.

\textbf{Damien Guard} drew the \href{https://damieng.com/zx-origins}{ZX
Origins} fonts --- hundreds of 8x8 faces, refined over thirty-five years
and given to anyone who wants them. A dozen of them are in the F7 font
menu. They are free to use and \emph{not} free to re-host, so this
machine does what it does with Commodore's chargen: it names them and
does not ship them. \verb|tools/mkzxfonts.py| turns your own download
into screen fonts; without it those menu entries fall back to the
default face, and nothing else breaks.

\textbf{Kenney} makes game art and puts it in the public domain, at a
scale and a standard that has quietly furnished a decade of small games.
The dungeon in the \verb|TINY| demo --- every tile, every little
person, and the map they stand on --- is his
\href{https://kenney.nl/assets/tiny-dungeon}{Tiny Dungeon}, CC0. The
demo exists because the art did. Two more games stand on the same
ground: \verb|SKYFIRE| flies Kenney's
\href{https://kenney.nl/assets/pixel-shmup}{Pixel Shmup} ships, and
\verb|FLUFFY| is drawn with \textbf{Chloe Wolfe}'s Game Boy platformer
set, also CC0. In both cases the art came first and the game was written
around it, which is the reverse of the usual order and a much better way
to spend an evening.

\textbf{Ian Schofield} wrote
\href{https://github.com/ijschofield/Tek40xx}{Tek40xx}, a Tektronix
4010/4014 storage tube on SDL2 that is also a telnet client. It rides
along on the machine's Linux as the second terminal
(Chapter~\ref{cha:linux}), built from upstream with one patch of ours.

\section{The bare metal, which the machine no longer stands on}
This section is kept because the debt is not cancelled by the code
coming out. From August to September 2026 the K4510 ran on a Raspberry
Pi~3B+ with no operating system underneath it, and that edition was
retired on 2026-09-07 (page~\pageref{sec:names}). Everything below is
why it was possible at all.

\textbf{Randy Rossi} wrote
\href{https://github.com/randyrossi/bmc64}{BMC64}: VICE on a Raspberry
Pi with no operating system under it, 50\,Hz smooth scrolling and input
latency measured in single frames. His initials were in that edition's
name --- the \textbf{BMC-K4510} was the K4510 on the bare-metal platform
he proved out --- and it is the direct reason the port existed at all:
the route to the Pi was always meant to be BMC64's \verb|emux_api| seam,
and the bare-metal case it proved is what made it thinkable. He also built the \textbf{VIC-II Kawari}, a
modern drop-in VIC-II with video modes the original never had: the
demonstration that you may extend an 8-bit machine's video chip and
still be working inside the tradition rather than outside it. VICKY is
a more reckless cousin of that idea. \textbf{minch} (aminch) has taken
BMC64 over from him and carries it onto the newer Pis; the project is
in good hands.

\textbf{Rene Stange} wrote \href{https://github.com/rsta2/circle}{Circle},
the bare-metal C++ environment that was the entire reason a Raspberry
Pi~3B+ could boot straight into this machine with no operating system
under it. \textbf{Xalior} wrote
\href{https://github.com/Xalior/circle-libsdl2}{circle-libsdl2}, which
let the desktop emulator move to the Pi essentially as written.

\section{The workshop}
The machine is built with \textbf{cc65} (compiler, assembler, linker
and runtime for the whole system ROM), Marco Baye's \textbf{ACME}
(Wozmon and the demos), Zsolt Soós's \textbf{64tass} (Tali Forth
and Supermon), \textbf{NASM}, \textbf{GCC} and \textbf{GNU Make}; it shows itself to you through
\textbf{SDL2} --- Sam Lantinga's library and its contributors', which
is the window, the keyboard and the sound on the desktop, and through
Xalior's port, on the Pi as well. This book is set in \textbf{XeLaTeX} with
\textbf{Clear Sans} and \textbf{Iosevka}, and every screenshot in it
was captured from the machine actually running --- a picture that
cannot be produced fails the build.

\section{Heritage}
Some debts are not code.

\textbf{Acorn Computers} --- the Tube, the sideways-ROM model, the
\verb|*| command prefix, and the idea that a second processor should
be a normal thing to own. The Beeb's ghost is all over this machine.

\textbf{Commodore} --- the other half of its soul: PETSCII, the SID,
and the line from the PET to the C65 that stopped one machine short of
this one.

\textbf{The MEGA65 project} --- for keeping the 45GS02 and the C65
dream alive in real silicon, and for open-roms.

\textbf{Kim Lemon} and the Lemoners --- \href{https://www.lemon64.com}{Lemon64}
has been the C64's front door since 1998: the games database, the
reviews and scans, the music, and a forum that actually answers.
Twenty-eight years of one person's dedication to a machine's community,
and a good half of the small facts this project needed came from there.

\textbf{The Meatloaf project} (idolpx) --- the C64 cartridge that made
a URL a file name, which is the whole of this machine's rule for the
network: name it, and it is fetched. \textbf{The FujiNet project} ---
the network device the Atari got first, whose N: device this machine
imitates and whose TNFS servers it speaks to;
\texttt{CD tnfs://fujinet.online/} is a directory on theirs.

\textbf{Paul Scott Robson} wrote the Neo6502's firmware --- the
operating system and BASIC that make an Olimex Neo6502 a computer
rather than a board --- and the handbook that documents it. That book
is the one this one is modelled on, down to the page size, and the rule
that every example must be run and photographed before it may be
printed is his. The other model is the \textbf{MEGA65 User's Guide}.

\textbf{The High Voltage SID Collection} and the composers in it, who
kept thirty years of a machine's music alive and free to listen to. The
K4510 played it for one summer, through a \verb|SIDPLAY| that retired
with the chips. None of that music was ever distributed with this
machine; it is theirs.

\begin{aside}
\textbf{The project's orchestrator} is Michael Longval --- Doc --- who
decided what the machine is, what it borrows and what it refuses, what
gets built next and what gets thrown away, and read every page of this
book against the machine.

\textbf{And the machine's other author.} Most of the K4510's own code
--- VICKY, SHEILA, the DMA engine, the ROM, K/OS, the Tube ULA, the
editors, RX --- was written in conversation with Claude
(Anthropic's Claude Code), over several months of long sessions with a
build log to prove it. The design decisions are the author's; a great
deal of the typing was not.
\end{aside}

\section{And whoever is missing}
This list was assembled by hand and is certainly incomplete. If your
work is in this machine and your name is not on this page, that is an
error and not a judgement --- open an issue at
\url{https://github.com/mlongval/k4510} and it will be fixed in the
next printing.
\endgroup
